\documentclass[12pt]{article}
\usepackage{amsmath, amssymb}
\usepackage[margin=1in]{geometry}
\setlength{\parskip}{6pt}
\title{QUBO Formulation: DNA Read Ordering Problem for Genome Assembly}
\date{}
\begin{document}
\maketitle

\subsection*{DNA Read Ordering Problem for Genome Assembly}

\paragraph{Problem Statement.}
Given a set of $N$ DNA reads indexed by $i \in \{0, \dots, N-1\}$, let $O_{ij}$ denote the overlap score between read $i$ and read $j$, defined as the length of the longest suffix of read $i$ that matches a prefix of read $j$. The goal is to find a linear ordering of all reads that maximizes the total overlap between consecutive reads in the sequence. Formally, we seek a permutation $\pi$ of $\{0, \dots, N-1\}$ that maximizes the objective function:
\begin{equation}
    Z(\pi) = \sum_{k=0}^{N-2} O_{\pi(k), \pi(k+1)}.
\end{equation}
This problem is equivalent to finding a Hamiltonian path in a complete directed graph with edge weights $O_{ij}$ that maximizes the total weight.

\paragraph{Decision Variables.}
We introduce binary variables $x_{i,k} \in \{0,1\}$ for all $i \in \{0, \dots, N-1\}$ and $k \in \{0, \dots, N-1\}$, where:
\begin{equation}
    x_{i,k} = \begin{cases}
        1 & \text{if read } i \text{ is placed at position } k \text{ in the ordering}, \\
        0 & \text{otherwise}.
    \end{cases}
\end{equation}
The variables form an $N \times N$ assignment matrix. A valid solution corresponds to a permutation matrix.

\paragraph{QUBO Derivation.}
We construct a Quadratic Unconstrained Binary Optimization (QUBO) formulation by converting the maximization problem into a minimization problem and incorporating constraints as quadratic penalty terms.

\textbf{Step 1: Objective Function.}
Maximizing $Z$ is equivalent to minimizing its negative. Substituting the variables $x_{i,k}$, the objective Hamiltonian is:
\begin{equation}
    H_{\text{obj}} = -\sum_{k=0}^{N-2} \sum_{i=0}^{N-1} \sum_{j=0}^{N-1} O_{ij} \, x_{i,k} \, x_{j,k+1}.
\end{equation}
The sum over $k$ runs from $0$ to $N-2$ because the last position $N-1$ has no successor in the path.

\textbf{Step 2: Constraints.}
The assignment must satisfy two equality constraints:
\begin{align}
    \sum_{k=0}^{N-1} x_{i,k} &= 1, \quad \forall i \in \{0, \dots, N-1\} \quad \text{(each read assigned once)}, \\
    \sum_{i=0}^{N-1} x_{i,k} &= 1, \quad \forall k \in \{0, \dots, N-1\} \quad \text{(each position filled once)}.
\end{align}

\textbf{Step 3: Penalty Terms.}
We introduce a penalty weight $\lambda > 0$ and convert constraints to quadratic penalties. Using the idempotent property $x_{i,k}^2 = x_{i,k}$ for binary variables, we expand the squared penalty terms.

For the first constraint (each read once), the penalty contribution for read $i$ is:
\begin{align}
    P_i^{(1)} &= \lambda \left( \sum_{k=0}^{N-1} x_{i,k} - 1 \right)^2 \\
    &= \lambda \left( \sum_{k=0}^{N-1} x_{i,k}^2 + \sum_{\substack{k,l=0 \\ k \neq l}}^{N-1} x_{i,k} x_{i,l} - 2 \sum_{k=0}^{N-1} x_{i,k} + 1 \right) \\
    &= \lambda \left( \sum_{k=0}^{N-1} x_{i,k} + \sum_{\substack{k,l=0 \\ k \neq l}}^{N-1} x_{i,k} x_{i,l} - 2 \sum_{k=0}^{N-1} x_{i,k} + 1 \right) \\
    &= \lambda \left( 1 - \sum_{k=0}^{N-1} x_{i,k} + \sum_{\substack{k,l=0 \\ k \neq l}}^{N-1} x_{i,k} x_{i,l} \right).
\end{align}
Summing over all reads $i$, the total penalty for the first constraint is:
\begin{equation}
    H_{\text{pen}}^{(1)} = \lambda \sum_{i=0}^{N-1} \left( 1 - \sum_{k=0}^{N-1} x_{i,k} + \sum_{\substack{k,l=0 \\ k \neq l}}^{N-1} x_{i,k} x_{i,l} \right).
\end{equation}

Similarly, for the second constraint (each position once), the penalty contribution for position $k$ is:
\begin{align}
    P_k^{(2)} &= \lambda \left( \sum_{i=0}^{N-1} x_{i,k} - 1 \right)^2 \\
    &= \lambda \left( 1 - \sum_{i=0}^{N-1} x_{i,k} + \sum_{\substack{i,j=0 \\ i \neq j}}^{N-1} x_{i,k} x_{j,k} \right).
\end{align}
Summing over all positions $k$, the total penalty for the second constraint is:
\begin{equation}
    H_{\text{pen}}^{(2)} = \lambda \sum_{k=0}^{N-1} \left( 1 - \sum_{i=0}^{N-1} x_{i,k} + \sum_{\substack{i,j=0 \\ i \neq j}}^{N-1} x_{i,k} x_{j,k} \right).
\end{equation}

\textbf{Step 4: Combined Hamiltonian.}
The total QUBO Hamiltonian $H(x)$ is the sum of the objective and penalty terms:
\begin{equation}
    H(x) = H_{\text{obj}} + H_{\text{pen}}^{(1)} + H_{\text{pen}}^{(2)}.
\end{equation}
Substituting the expressions:
\begin{align}
    H(x) &= -\sum_{k=0}^{N-2} \sum_{i,j} O_{ij} x_{i,k} x_{j,k+1} \nonumber \\
    &\quad + \lambda \sum_{i} \left( 1 - \sum_{k} x_{i,k} + \sum_{\substack{k,l \\ k \neq l}} x_{i,k} x_{i,l} \right) \nonumber \\
    &\quad + \lambda \sum_{k} \left( 1 - \sum_{i} x_{i,k} + \sum_{\substack{i,j \\ i \neq j}} x_{i,k} x_{j,k} \right).
\end{align}

\textbf{Step 5: Q Matrix Entries and Offset.}
We express $H(x)$ in the standard QUBO form $H(x) = \sum_{u,v} Q_{u,v} x_u x_v + c$, where the variables are indexed by a single index $u = iN+k$. The constant offset $c$ collects all terms independent of $x$:
\begin{equation}
    c = \lambda \sum_{i=0}^{N-1} 1 + \lambda \sum_{k=0}^{N-1} 1 = N\lambda + N\lambda = 2N\lambda.
\end{equation}
The linear terms $x_{i,k}$ arise from $-\lambda \sum_i \sum_k x_{i,k}$ and $-\lambda \sum_k \sum_i x_{i,k}$, yielding a coefficient of $-2\lambda$ for each $x_{i,k}$. Thus, the diagonal entries of $Q$ are:
\begin{equation}
    Q_{(i,k),(i,k)} = -2\lambda, \quad \forall i,k.
\end{equation}
The quadratic terms $x_{i,k} x_{j,l}$ with $(i,k) \neq (j,l)$ arise from the penalty expansions and the objective. Collecting coefficients:
\begin{align}
    Q_{(i,k),(i,l)} &= \lambda, \quad \forall i, \, \forall k \neq l, \\
    Q_{(i,k),(j,k)} &= \lambda, \quad \forall k, \, \forall i \neq j, \\
    Q_{(i,k),(j,k+1)} &= -O_{ij}, \quad \forall k \in \{0, \dots, N-2\}, \, \forall i,j.
\end{align}
All other entries $Q_{(i,k),(j,l)}$ are zero. The matrix $Q$ is symmetric, so $Q_{(j,k+1),(i,k)} = -O_{ij}$.

\paragraph{Penalty Weight Justification.}
Let $W = \max_{i,j} |O_{ij}|$ be the maximum absolute overlap score. We require $\lambda > W$. This choice ensures that the penalty cost for violating any constraint strictly exceeds the maximum possible gain in the objective function from any single adjacency term. Specifically, the minimum energy increase due to a constraint violation is $\lambda$, while the maximum reduction in the objective (increase in overlap) from any edge is bounded by $W$. Thus, $\lambda > W$ guarantees that the global minimum of $H(x)$ cannot occur at an infeasible point, as any infeasible solution would incur a penalty strictly greater than the best possible objective improvement.

\paragraph{Correctness Argument.}
Minimizing $H(x)$ over $x \in \{0,1\}^{N^2}$ recovers the optimal read ordering. For any feasible assignment $x$, the penalty terms vanish, and $H(x) = H_{\text{obj}}(x) = -Z(x)$. Thus, minimizing $H$ over feasible solutions is equivalent to maximizing the total overlap $Z$. For any infeasible $x$, at least one constraint is violated, resulting in a penalty term of at least $\lambda$. Since $\lambda > W$ and the maximum possible objective value is bounded, the energy of any infeasible solution is strictly greater than that of the optimal feasible solution. Consequently, the global minimum of $H(x)$ corresponds to a feasible permutation that maximizes the total overlap.

\paragraph{Illustrative Example.}
Consider an instance with $N=3$ reads $r_0, r_1, r_2$. Suppose the overlap matrix $O$ is given by:
\begin{equation}
    O = \begin{pmatrix}
        0 & 2 & 0 \\
        0 & 0 & 2 \\
        0 & 0 & 0
    \end{pmatrix},
\end{equation}
where $O_{0,1}=2$ and $O_{1,2}=2$, indicating overlaps $r_0 \to r_1$ and $r_1 \to r_2$. Let $\lambda = 3$, satisfying $\lambda > \max |O_{ij}| = 2$.

The offset is $c = 2N\lambda = 6 \times 3 = 18$.
The diagonal entries are $Q_{(i,k),(i,k)} = -2\lambda = -6$.
The penalty off-diagonals are $Q_{(i,k),(i,l)} = 3$ for $k \neq l$ and $Q_{(i,k),(j,k)} = 3$ for $i \neq j$.
The objective terms are $Q_{(0,0),(1,1)} = -O_{0,1} = -2$ and $Q_{(1,1),(2,2)} = -O_{1,2} = -2$.

The Hamiltonian includes terms such as:
\begin{align}
    H(x) &= 18 - 6 \sum_{i,k} x_{i,k} + 3 \sum_{i} \sum_{k \neq l} x_{i,k} x_{i,l} + 3 \sum_{k} \sum_{i \neq j} x_{i,k} x_{j,k} \nonumber \\
    &\quad - 2 x_{0,0} x_{1,1} - 2 x_{1,1} x_{2,2}.
\end{align}
The optimal solution corresponds to the ordering $r_0 \to r_1 \to r_2$, represented by $x_{0,0}=1, x_{1,1}=1, x_{2,2}=1$ (all other $x_{i,k}=0$). The energy is:
\begin{equation}
    H(x^*) = 18 - 6(3) + 0 + 0 - 2(1) - 2(1) = 18 - 18 - 4 = -4.
\end{equation}
This matches the negative of the total overlap $Z = 2+2=4$. Any infeasible assignment would yield $H(x) \ge 3$, confirming optimality.

\end{document}